% onlymath.tex -- standalone mathematical companion to sample.tex.
%
% Every formula in sample.tex, pulled out with the derivation shown
% (not just the final restated equation). Single-column article class
% since this is a technical appendix/study note, not a submission
% draft -- pair with sample.tex for the full paper.
%
% Compiled and verified with tectonic in this session; zero errors.

\documentclass[11pt]{article}
\usepackage{amsmath,amssymb}
\usepackage[T1]{fontenc}
\usepackage{geometry}
\geometry{margin=1in}

\title{Fuzzy Phonemes: Mathematical Derivations}
\author{Companion to \texttt{sample.tex}}
\date{}

\begin{document}
\maketitle

\section{Candidate selection as Bayesian posterior}

\subsection{Setup}
Given an ASR hypothesis word $w_{asr}$, a candidate set
$\mathcal{C} = \{c_1, \ldots, c_k\}$, the acoustic segment $s$, and
sentence context $T$, we want the posterior probability of each
candidate word $w \in \mathcal{C} \cup \{w_{asr}\}$ being the one
actually spoken.

\subsection{Derivation}
Start from Bayes' rule over the joint evidence $(s, T)$:
\begin{equation}
P(w \mid s, T) = \frac{P(s, T \mid w) \, P(w)}{P(s, T)}.
\end{equation}
Assume the acoustic evidence $s$ and the contextual evidence $T$ are
conditionally independent given the word $w$ (a standard naive-Bayes-style
simplifying assumption --- acoustic realization and surrounding topic are
treated as separate, non-interacting evidence channels):
\begin{equation}
P(s, T \mid w) = P(s \mid w) \, P(T \mid w).
\end{equation}
Parameterize each factor: $P(s \mid w) = P(s \mid w, \theta_{ac})$ (the
acoustic model, itself a fusion of general acoustic features, explicit
phonetic cues, and the GP classifier's output), $P(T \mid w) = P(T \mid
w, \theta_{ctx})$ (the LDA context model), and the prior $P(w) = P(w \mid
\theta_{pr})$ (corpus word-frequency prior). Substituting:
\begin{equation}
P(w \mid s, T, \theta) \propto P(s \mid w, \theta_{ac}) \cdot
P(T \mid w, \theta_{ctx}) \cdot P(w \mid \theta_{pr}).
\end{equation}
The proportionality constant $P(s,T)$ doesn't depend on $w$, so it's
recovered by normalizing over the candidate set --- summing the same
unnormalized expression over every $w' \in \mathcal{C} \cup \{w_{asr}\}$
and dividing:
\begin{equation}
P(w \mid s, T, \theta) = \frac{P(s \mid w, \theta_{ac}) \cdot P(T \mid w, \theta_{ctx}) \cdot P(w \mid \theta_{pr})}
{\displaystyle\sum_{w' \in \mathcal{C} \cup \{w_{asr}\}} P(s \mid w', \theta_{ac}) \cdot P(T \mid w', \theta_{ctx}) \cdot P(w' \mid \theta_{pr})}.
\label{eq:posterior}
\end{equation}
The decision rule is $w^* = \arg\max_{w \in \mathcal{C}\cup\{w_{asr}\}} P(w \mid s, T, \theta)$.

\section{Gaussian Process classification}

\subsection{Feature vector}
Each candidate occurrence is represented by:
\begin{equation}
X = [\text{VOT},\ F0_{\text{pert}},\ H1H2,\ \text{AspDur},\ \text{GMM},\ \text{DTW},\ \text{Fuzzy},\ \text{Topic}] \in \mathbb{R}^8.
\end{equation}
Standardized per-dimension using training-split statistics only:
\begin{equation}
X'_j = \frac{X_j - \mu_j^{\text{train}}}{\sigma_j^{\text{train}}}, \qquad j = 1, \ldots, 8.
\label{eq:standardize}
\end{equation}

\subsection{GP prior and the classification link}
A GP places a prior directly over a latent function $f(X)$:
\begin{equation}
f(X) \sim \mathcal{GP}\big(0,\ k(X, X')\big),
\end{equation}
i.e.\ any finite set of function values is jointly Gaussian, fully
specified by the covariance (kernel) function $k$. For \emph{regression}
this $f$ would be the target directly; for \emph{classification}, binary
labels $y \in \{0,1\}$ can't be Gaussian, so $f$ is squashed through a
sigmoid-like link before it becomes a probability. This model uses the
\textbf{probit} link (the standard normal CDF $\Phi$) with a Bernoulli
likelihood:
\begin{align}
p(y=1 \mid f) &= \Phi(f), \\
p(y \mid f) &= \Phi(f)^y \, \big(1-\Phi(f)\big)^{1-y}.
\end{align}
Because the Bernoulli likelihood isn't Gaussian, the true posterior
$p(f \mid \mathcal{D})$ has no closed form. \textbf{Expectation
Propagation (EP)} approximates it by iteratively replacing each
non-Gaussian likelihood factor with a Gaussian approximation, matching
moments (mean and variance) against the true tilted distribution at each
step, until the approximation converges. The predictive distribution for
a new point $X_*$ is then obtained by pushing this approximate Gaussian
posterior over $f_*$ through $\Phi(\cdot)$ again to get a class
probability.

\subsection{Covariance functions compared}
Squared Exponential (RBF) --- infinitely differentiable, assumes very
smooth $f$:
\begin{equation}
k_{\text{RBF}}(x,x') = \sigma_f^2 \exp\left(-\frac{\|x-x'\|^2}{2\ell^2}\right).
\end{equation}
Mat\'ern with $\nu = 5/2$ --- twice mean-square differentiable, a milder
smoothness assumption obtained by taking the general Mat\'ern kernel
\begin{equation}
k_{\text{M}\nu}(x,x') = \sigma_f^2 \frac{2^{1-\nu}}{\Gamma(\nu)}
\left(\frac{\sqrt{2\nu}\|x-x'\|}{\ell}\right)^{\nu}
K_\nu\!\left(\frac{\sqrt{2\nu}\|x-x'\|}{\ell}\right)
\end{equation}
and substituting the half-integer case $\nu=5/2$, where the modified
Bessel function $K_\nu$ collapses to an elementary closed form:
\begin{equation}
k_{\text{M52}}(x,x') = \sigma_f^2\left(1+\frac{\sqrt{5}\|x-x'\|}{\ell}+\frac{5\|x-x'\|^2}{3\ell^2}\right)\exp\left(-\frac{\sqrt{5}\|x-x'\|}{\ell}\right).
\end{equation}
In both, $\sigma_f^2$ is the signal variance and $\ell$ the lengthscale;
both are hyperparameters fit by maximizing the log marginal likelihood
$\log p(\mathbf{y} \mid X)$ via L-BFGS-B (a quasi-Newton method using an
approximate inverse Hessian, avoiding the cost of computing the true
Hessian at every step).

\subsection{Regression-to-classification}
A second arm trains a GP \emph{regressor} (identical kernel choices, but
a Gaussian likelihood instead of Bernoulli/probit) on a continuous
confusability score $y \in [0,1]$, then thresholds the predictive mean:
\begin{equation}
\hat{y}_{\text{class}} =
\begin{cases}
1 & \hat{y}_{\text{reg}} \geq \tau \\
0 & \hat{y}_{\text{reg}} < \tau
\end{cases}
\qquad \tau = 0.5.
\label{eq:threshold}
\end{equation}
\emph{Open question (unresolved in the source proposal):} the definition
of $y$ itself. The working assumption used elsewhere in this project is
$y=1$ for the true candidate and $y=0$ otherwise --- i.e.\ the regressor
is fit on the same labels the classifier uses, just modeled continuously.

\section{Fuzzy evidence aggregation}

For a cue value $c$ (e.g.\ VOT in ms), define membership functions
$\mu_m(c) \in [0,1]$ for each linguistic category $m$ (e.g.\ Unaspirated,
Aspirated, Ambiguous), fit from the training-data distribution of that
cue. The aggregated fuzzy score for a candidate is a weighted sum over
all cues $c^{(1)}, \ldots, c^{(n)}$ and their most-supported memberships:
\begin{equation}
\text{Fuzzy}(w) = \sum_{i=1}^{n} \beta_i \cdot \max_m \, \mu_m\!\left(c^{(i)}\right),
\end{equation}
where the weights $\beta_i$ are learned during training (the exact
fitting procedure is an open question --- see \texttt{sushank.md} \S2.3
--- one reasonable choice is fitting $\beta$ by logistic regression
against the same training labels used elsewhere).

\section{Contextual evidence via LDA}

\subsection{Generative model (brief)}
LDA assumes each document is a mixture over $K$ latent topics, and each
topic is a distribution over the vocabulary. For a document, a topic
proportion vector $\theta \sim \text{Dirichlet}(\alpha)$ is drawn, then
each word's topic assignment $z \sim \text{Categorical}(\theta)$, then
the word itself $\sim \text{Categorical}(\phi_z)$, where $\phi_z \sim
\text{Dirichlet}(\eta)$ is topic $z$'s word distribution. Fitting via
\texttt{gensim} estimates $\phi$ (topic-word distributions) from the
50{,}000-document corpus; $K=20$.

\subsection{Context score}
For a candidate word $w$ and its surrounding sentence context (excluding
$w$ itself), the posterior topic distribution $P(t \mid \text{context})$
is inferred, then combined with each topic's affinity for $w$:
\begin{equation}
\text{ContextScore}(w) = \sum_{t=1}^{K} P(t \mid \text{context}) \cdot P(w \mid t).
\label{eq:context}
\end{equation}
This is precisely $\mathbb{E}_{t \sim P(\cdot \mid \text{context})}\big[P(w \mid t)\big]$
--- the expected topic-word probability of $w$ under the context's
inferred topic mixture. A candidate that fits the surrounding topic well
gets a high score even if its acoustic evidence alone is ambiguous.

\section{Evaluation metrics}

These are standard definitions (not stated as explicit equations in the
source proposal, but implied by naming the metrics) --- included here so
the evaluation code has an unambiguous target to implement against.

\subsection{Brier score}
For $N$ predictions with probability $p_i$ assigned to the outcome that
actually occurred $o_i \in \{0,1\}$:
\begin{equation}
\text{Brier} = \frac{1}{N} \sum_{i=1}^{N} (p_i - o_i)^2.
\end{equation}
Lower is better; a perfectly calibrated, always-correct classifier
scores 0.

\subsection{Expected Calibration Error (ECE)}
Partition predictions into $B$ confidence bins by predicted probability.
For bin $b$, let $\text{acc}(b)$ be the empirical accuracy of predictions
in that bin and $\text{conf}(b)$ the average predicted confidence:
\begin{equation}
\text{ECE} = \sum_{b=1}^{B} \frac{|b|}{N} \, \big|\text{acc}(b) - \text{conf}(b)\big|,
\end{equation}
where $|b|$ is the number of predictions falling in bin $b$. ECE is the
confidence-weighted average gap between how confident the model claims
to be and how often it's actually right.

\end{document}
